\theory{Dimension theory}{How can the size or number of independent directions of an algebraic object be measured when the object is not a simple line, plane, or vector space?}{Dimension theory studies several notions of dimension and the conditions under which they agree. In commutative algebra and algebraic geometry, the main notion is Krull dimension: the maximum length of a chain of prime ideals. It measures how many layers of algebraic subspaces a ring or scheme can contain.}{Algebraic geometry, commutative algebra, singularities, algebraic number theory, module theory, homological algebra, and the comparison of geometric and algebraic complexity.}{Prime chains, local rings, and homological dimensions turn the informal size of an object into computable algebraic quantities.}

\paragraph{Wolfgang Krull}
Krull defined the dimension of a commutative ring through chains of prime ideals. This algebraic definition measures how many successive layers of algebraic subspaces can occur.

\paragraph{Henri Lebesgue and topological dimension theorists}
Lebesgue and later covering-dimension researchers developed ways to assign dimension to general topological spaces. These notions allow algebraic, geometric, and topological dimensions to be compared when their hypotheses overlap.
\end{historylist}
\section*{3. Theory}
\subsection*{Core theory}
\paragraph{Why dimension needs a theory}

For a vector space, dimension is the number of vectors in a basis. For a line, plane, or manifold, it is the number of independent local coordinates. Algebraic sets can be singular, reducible, or infinitely complicated, so these familiar definitions may no longer agree or may not even apply.

Dimension theory studies several questions. How long can a chain of proper algebraic subspaces be? How many independent parameters describe a family? How many equations are needed locally? How complicated is a module's resolution? For regular geometric objects these answers often coincide; for singular or non-Noetherian objects they may differ. \cite{dimension_theory_ref1,dimension_theory_ref2}

\paragraph{Krull dimension and prime chains}

Let $R$ be a commutative ring. A prime ideal is a set of elements closed under addition and multiplication by arbitrary ring elements such that if a product $ab$ lies in the ideal, then $a$ or $b$ lies in it. Prime ideals represent irreducible algebraic conditions.

The Krull dimension of $R$ is the supremum of the lengths of chains
\[
 \mathfrak p_0\subsetneq\mathfrak p_1\subsetneq\cdots\subsetneq\mathfrak p_n
\]
of prime ideals. The height of a prime $\mathfrak p$ is the dimension of the localization $R_{\mathfrak p}$, equivalently the maximum length of a chain ending at $\mathfrak p$.

For a field $k$, the ring $k$ has dimension zero because its only prime ideal is $(0)$. The polynomial ring $k[x]$ has the chain $(0)\subsetneq(x)$ and dimension one. The ring $k[x,y]$ has
\[
 (0)\subsetneq(x)\subsetneq(x,y),
\]
so its dimension is two. Algebraically, the plane has two independent coordinates and two layers of prime conditions.

\paragraph{Polynomial rings, varieties, and chains}

For a Noetherian ring $R$, adjoining one polynomial variable raises the dimension:
\[
 \dim R[x]=\dim R+1.
\]
Thus a polynomial ring in $n$ variables over a field has dimension $n$. Quotients lower dimension according to the equations imposed. A hypersurface defined by one nonzero equation in affine $n$-space usually has dimension $n-1$, although singularities and components require care.

The prime spectrum $\operatorname{Spec}R$ turns the set of prime ideals into a geometric space. Inclusion of primes becomes specialization: a larger prime represents a more constrained point. Krull dimension is then the dimension of the longest chain of geometric strata. Localization focuses on a neighborhood, while quotienting by a prime restricts to a closed irreducible subspace.

For finitely generated algebras over a field, algebraic, geometric, and transcendence dimensions agree in the expected cases. The dimension of an irreducible variety is the transcendence degree of its function field over the ground field. \cite{dimension_theory_ref1,dimension_theory_ref3}

\paragraph{Noetherian rings and principal ideal bounds}

A ring is Noetherian when every ascending chain of ideals eventually stops, or equivalently when every ideal has a finite generating set. This condition makes dimension theory tractable. Polynomial rings over Noetherian rings remain Noetherian by the Hilbert basis theorem.

Krull's principal ideal theorem says that a minimal prime over a principal ideal generated by a nonzero nonunit has height at most one. More generally, a minimal prime over an ideal generated by $r$ elements has height at most $r$ under standard Noetherian hypotheses. This formalizes the geometric expectation that one equation cuts dimension by at most one.

A ring is catenary when saturated prime chains between two fixed primes have the same length. Catenarity ensures that dimension behaves predictably under passage between strata. Many finitely generated algebras over fields are catenary, while arbitrary non-Noetherian rings may violate familiar dimension formulas. \cite{dimension_theory_ref2,dimension_theory_ref4}

\paragraph{Regularity, singularities, and local dimension}

At a point of an algebraic variety, the local ring records functions defined near that point. A regular local ring has dimension equal to the dimension of its cotangent or tangent space. If the tangent space is larger than the local Krull dimension, the point is singular: too many first-order directions meet there.

For example, a crossing such as $xy=0$ has two branches meeting at the origin. The local ring has dimension one, but the tangent space has two independent directions, signaling a singularity. A smooth curve has matching dimensions and tangent directions.

Regular rings are important because many dimension notions agree there. Cohen--Macaulay and Gorenstein conditions weaken or refine regularity by controlling depth, duality, and resolutions. These properties help determine whether intersections have the expected dimension and whether algebraic invariants behave like those of a smooth space. \cite{dimension_theory_ref5}

\paragraph{Homological and module dimensions}

Krull dimension measures prime-ideal chains. Homological dimension measures how many steps are needed to resolve a module by free or projective modules. The global dimension of a ring is the supremum of projective dimensions of its modules. For a regular Noetherian ring, the global dimension agrees with the Krull dimension in important settings.

Depth measures the length of a regular sequence in a local ring. The inequality
\[
 \operatorname{depth}R\leq\dim R
\]
is fundamental. Equality characterizes Cohen--Macaulay local rings. Thus dimension theory connects the geometry of strata with the algebra of equations and the complexity of solving them.

For modules, dimension can also mean length, rank, or support dimension. A module supported only at finitely many points has dimension zero, while a module whose support fills a surface has dimension two. The correct notion depends on the question being asked.

\paragraph{Valuation, non-Noetherian, and noncommutative dimensions}

A valuation ring has ideals totally ordered by inclusion. Its prime spectrum is correspondingly simple, and its dimension can often be read from the chain of prime ideals. Valuation rings provide local models in algebraic geometry and number theory.

For non-Noetherian rings, ideals need not be finitely generated, chains may be very long, and polynomial dimension formulas can fail. Little is known in full generality, so the theory studies special classes such as coherent, Prüfer, or valuative domains.

Noncommutative rings do not have prime spectra with all the same geometric behavior. One can define Krull dimension-like invariants using chains of prime ideals or modules, but left and right dimensions may differ. These notions are used in representation theory and the study of operator algebras. \cite{dimension_theory_ref6}

\paragraph{What the theory depends on and where it is useful}

Dimension theory depends on rings, ideals, modules, localization, polynomial algebra, algebraic geometry, and homological algebra. It helps algebraic geometry by measuring varieties and fibers, number theory by measuring spectra of rings of integers, and singularity theory by comparing tangent and local dimensions.

It is useful whenever a problem has several possible notions of size. A theorem that equates two dimensions is valuable because it allows a geometric calculation to be replaced by an algebraic one. A mismatch is equally informative: it can reveal a singularity, a bad intersection, or a non-Noetherian pathology.

\section*{4. Important theorems and results}
\begin{enumerate}[leftmargin=*,itemsep=0.35em]

\item \textbf{Polynomial-dimension theorem.} For a Noetherian ring or valuation ring, $\dim R[x]=\dim R+1$ under the standard hypotheses. \cite{dimension_theory_ref1}

\item \textbf{Hilbert basis theorem.} If $R$ is Noetherian, then $R[x]$ is Noetherian. This is a main reason polynomial dimension can be studied inductively. \cite{dimension_theory_ref2}

\item \textbf{Krull principal ideal theorem.} A minimal prime over a principal ideal has height at most one; the version for $r$ generators gives height at most $r$ under Noetherian assumptions. \cite{dimension_theory_ref4}

\item \textbf{Dimension--transcendence theorem.} For a finitely generated domain over a field, the dimension equals the transcendence degree of its fraction field over the base field. \cite{dimension_theory_ref3}

\item \textbf{Regularity criterion.} A Noetherian local ring is regular when its Krull dimension equals the dimension of its maximal ideal modulo its square. \cite{dimension_theory_ref5}
\end{enumerate}

\theoryapplications
\theoryconnections{dimension theory}{%
\item Commutative algebra supplies prime ideals and chains, algebraic geometry interprets them as geometric dimension, and homological algebra supplies projective and injective dimensions.
\item Noetherian hypotheses are important because they prevent uncontrolled infinite chains.
}{%
\item Dimension theory helps algebraic geometry compare varieties, helps local algebra measure singularities, and helps homological algebra predict the length of resolutions.
\item A dimension is useful only after its definition and category are stated; different dimensions answer different questions.
}
\section*{7. Sample Questions}
\emph{Exercise reference:} \cite{dimension_theory_ref1}. The cited edition is the source for the exercise set; consult its Exercises section for the official statement and numbering.
\begin{enumerate}[leftmargin=*,itemsep=0.9em]

\item \textbf{Exercise 1.} What is the Krull dimension of $k[x]$?

\emph{Solution.} The prime chain $(0)\subsetneq(x)$ has length one, and no longer chain exists. Hence $\dim k[x]=1$.

\item \textbf{Exercise 2.} Give a prime chain showing that $k[x,y]$ has dimension at least two.

\emph{Solution.} The ideals $(0)$, $(x)$, and $(x,y)$ are prime and strictly nested. The chain has length two, so the dimension is at least two; the polynomial-dimension theorem shows it is exactly two.

\item \textbf{Exercise 3.} What does height measure?

\emph{Solution.} The height of a prime ideal is the maximum number of strict prime inclusions in a chain ending at that ideal. It measures how many algebraic layers lie below the corresponding subspace.

\item \textbf{Exercise 4.} Why does the crossing $xy=0$ signal a singularity at the origin?

\emph{Solution.} The curve has local dimension one, but its two branches give two independent tangent directions at the origin. The tangent dimension is therefore larger than the local dimension, which is the regularity test for singularity.

\item \textbf{Exercise 5.} Why are Noetherian hypotheses important?

\emph{Solution.} They ensure finite generation and prevent endlessly increasing chains of ideals. Without them, polynomial dimension formulas, finite-height arguments, and geometric interpretations can fail.

\end{enumerate}
\section*{8. References}


\begin{thebibliography}{9}
\bibitem{dimension_theory_ref1} H. Matsumura, \emph{Commutative Ring Theory}, Cambridge University Press, 1986.
\bibitem{dimension_theory_ref2} M. F. Atiyah and I. G. Macdonald, \emph{Introduction to Commutative Algebra}, Addison-Wesley, 1969.
\bibitem{dimension_theory_ref3} R. Hartshorne, \emph{Algebraic Geometry}, Springer, 1977.
\bibitem{dimension_theory_ref4} W. Krull, "Idealtheorie in Ringen ohne Endlichkeitsbedingung," \emph{Mathematische Annalen}, 1926.
\bibitem{dimension_theory_ref5} H. Matsumura, \emph{Commutative Algebra}, Benjamin, 1980.
\bibitem{dimension_theory_ref6} J. C. McConnell and J. C. Robson, \emph{Noncommutative Noetherian Rings}, Wiley, 1987.
\end{thebibliography}
